\documentclass[12pt]{article}

% Idioma y codificación
\usepackage[spanish, es-tabla]{babel}       %es-tabla para que se titule "Tabla"
\usepackage[utf8]{inputenc}

% Márgenes
\usepackage[a4paper,top=3cm,bottom=2.5cm,left=3cm,right=3cm]{geometry}

% Comentarios de bloque
\usepackage{verbatim}

% Paquetes de links
\usepackage[hidelinks]{hyperref}    % Permite enlaces
\usepackage{url}                    % redirecciona a la web

% Más opciones para enumeraciones
\usepackage{enumitem}

% Personalizar la portada
\usepackage{titling}


% Paquetes de tablas
\usepackage{multirow}


%------------------------------------------------------------------------

%Paquetes de figuras
\usepackage{caption}
\usepackage{subcaption} % Figuras al lado de otras
\usepackage{float}      % Poner figuras en el sitio indicado H.


% Paquetes de imágenes
\usepackage{graphicx}       % Paquete para añadir imágenes
\usepackage{transparent}    % Para manejar la opacidad de las figuras

% Paquete para usar colores
\usepackage[dvipsnames]{xcolor}
\usepackage{pagecolor}      % Para cambiar el color de la página

% Habilita tamaños de fuente mayores
\usepackage{fix-cm}


%------------------------------------------------------------------------

% Paquetes de matemáticas
\usepackage{mathtools, amsfonts, amssymb, mathrsfs}
\usepackage[makeroom]{cancel}     % Simplificar tachando
\usepackage{polynom}    % Divisiones y Ruffini
\usepackage{units} % Para poner fracciones diagonales con \nicefrac

\usepackage{pgfplots}   %Representar funciones
\pgfplotsset{compat=1.18}  % Versión 1.18

\usepackage{tikz-cd}    % Para usar diagramas de composiciones
\usetikzlibrary{calc}   % Para usar cálculo de coordenadas en tikz

%Definición de teoremas, etc.
\usepackage{amsthm}
%\swapnumbers   % Intercambia la posición del texto y de la numeración

\theoremstyle{plain}

\makeatletter
\@ifclassloaded{article}{
  \newtheorem{teo}{Teorema}[section]
}{
  \newtheorem{teo}{Teorema}[chapter]  % Se resetea en cada chapter
}
\makeatother

\newtheorem{coro}{Corolario}[teo]           % Se resetea en cada teorema
\newtheorem{prop}[teo]{Proposición}         % Usa el mismo contador que teorema
\newtheorem{lema}[teo]{Lema}                % Usa el mismo contador que teorema

\theoremstyle{remark}
\newtheorem*{observacion}{Observación}

\theoremstyle{definition}

\makeatletter
\@ifclassloaded{article}{
  \newtheorem{definicion}{Definición} [section]     % Se resetea en cada chapter
}{
  \newtheorem{definicion}{Definición} [chapter]     % Se resetea en cada chapter
}
\makeatother

\newtheorem*{notacion}{Notación}
\newtheorem*{ejemplo}{Ejemplo}
\newtheorem*{ejercicio*}{Ejercicio}             % No numerado
\newtheorem{ejercicio}{Ejercicio} [section]     % Se resetea en cada section


% Modificar el formato de la numeración del teorema "ejercicio"
\renewcommand{\theejercicio}{%
  \ifnum\value{section}=0 % Si no se ha iniciado ninguna sección
    \arabic{ejercicio}% Solo mostrar el número de ejercicio
  \else
    \thesection.\arabic{ejercicio}% Mostrar número de sección y número de ejercicio
  \fi
}


% \renewcommand\qedsymbol{$\blacksquare$}         % Cambiar símbolo QED
%------------------------------------------------------------------------

% Paquetes para encabezados
\usepackage{fancyhdr}
\pagestyle{fancy}
\fancyhf{}

\newcommand{\helv}{ % Modificación tamaño de letra
\fontfamily{}\fontsize{12}{12}\selectfont}
\setlength{\headheight}{15pt} % Amplía el tamaño del índice


%\usepackage{lastpage}   % Referenciar última pag   \pageref{LastPage}
\fancyfoot[C]{\thepage}

%------------------------------------------------------------------------

% Conseguir que no ponga "Capítulo 1". Sino solo "1."
\makeatletter
\@ifclassloaded{book}{
  \renewcommand{\chaptermark}[1]{\markboth{\thechapter.\ #1}{}} % En el encabezado
    
  \renewcommand{\@makechapterhead}[1]{%
  \vspace*{50\p@}%
  {\parindent \z@ \raggedright \normalfont
    \ifnum \c@secnumdepth >\m@ne
      \huge\bfseries \thechapter.\hspace{1em}\ignorespaces
    \fi
    \interlinepenalty\@M
    \Huge \bfseries #1\par\nobreak
    \vskip 40\p@
  }}
}
\makeatother

%------------------------------------------------------------------------
% Paquetes de cógido
\usepackage{minted}
\renewcommand\listingscaption{Código fuente}

\usepackage{fancyvrb}
% Personaliza el tamaño de los números de línea
\renewcommand{\theFancyVerbLine}{\small\arabic{FancyVerbLine}}

% Estilo para C++
\newminted{cpp}{
    frame=lines,
    framesep=2mm,
    baselinestretch=1.2,
    linenos,
    escapeinside=||
}



\usepackage{listings} % Para incluir código desde un archivo

\renewcommand\lstlistingname{Código Fuente}
\renewcommand\lstlistlistingname{Índice de Códigos Fuente}

% Definir colores
\definecolor{vscodepurple}{rgb}{0.5,0,0.5}
\definecolor{vscodeblue}{rgb}{0,0,0.8}
\definecolor{vscodegreen}{rgb}{0,0.5,0}
\definecolor{vscodegray}{rgb}{0.5,0.5,0.5}
\definecolor{vscodebackground}{rgb}{0.97,0.97,0.97}
\definecolor{vscodelightgray}{rgb}{0.9,0.9,0.9}

% Configuración para el estilo de C similar a VSCode
\lstdefinestyle{vscode_C}{
  backgroundcolor=\color{vscodebackground},
  commentstyle=\color{vscodegreen},
  keywordstyle=\color{vscodeblue},
  numberstyle=\tiny\color{vscodegray},
  stringstyle=\color{vscodepurple},
  basicstyle=\scriptsize\ttfamily,
  breakatwhitespace=false,
  breaklines=true,
  captionpos=b,
  keepspaces=true,
  numbers=left,
  numbersep=5pt,
  showspaces=false,
  showstringspaces=false,
  showtabs=false,
  tabsize=2,
  frame=tb,
  framerule=0pt,
  aboveskip=10pt,
  belowskip=10pt,
  xleftmargin=10pt,
  xrightmargin=10pt,
  framexleftmargin=10pt,
  framexrightmargin=10pt,
  framesep=0pt,
  rulecolor=\color{vscodelightgray},
  backgroundcolor=\color{vscodebackground},
}

%------------------------------------------------------------------------

% Comandos definidos
\newcommand{\bb}[1]{\mathbb{#1}}
\newcommand{\cc}[1]{\mathcal{#1}}

% I prefer the slanted \leq
\let\oldleq\leq % save them in case they're every wanted
\let\oldgeq\geq
\renewcommand{\leq}{\leqslant}
\renewcommand{\geq}{\geqslant}

% Si y solo si
\newcommand{\sii}{\iff}

% Letras griegas
\newcommand{\eps}{\epsilon}
\newcommand{\veps}{\varepsilon}
\newcommand{\lm}{\lambda}

\newcommand{\ol}{\overline}
\newcommand{\ul}{\underline}
\newcommand{\wt}{\widetilde}
\newcommand{\wh}{\widehat}

\let\oldvec\vec
\renewcommand{\vec}{\overrightarrow}

% Derivadas parciales
\newcommand{\del}[2]{\frac{\partial #1}{\partial #2}}
\newcommand{\Del}[3]{\frac{\partial^{#1} #2}{\partial^{#1} #3}}
\newcommand{\deld}[2]{\dfrac{\partial #1}{\partial #2}}
\newcommand{\Deld}[3]{\dfrac{\partial^{#1} #2}{\partial^{#1} #3}}


\newcommand{\AstIg}{\stackrel{(\ast)}{=}}
\newcommand{\Hop}{\stackrel{L'H\hat{o}pital}{=}}

\newcommand{\red}[1]{{\color{red}#1}} % Para integrales, destacar los cambios.

% Método de integración
\newcommand{\MetInt}[2]{
    \left[\begin{array}{c}
        #1 \\ #2
    \end{array}\right]
}

% Declarar aplicaciones
% 1. Nombre aplicación
% 2. Dominio
% 3. Codominio
% 4. Variable
% 5. Imagen de la variable
\newcommand{\Func}[5]{
    \begin{equation*}
        \begin{array}{rrll}
            #1:& #2 & \longrightarrow & #3\\
               & #4 & \longmapsto & #5
        \end{array}
    \end{equation*}
}

%------------------------------------------------------------------------

\usepackage{pgfplots}
\pgfplotsset{compat=1.15}
\usepackage{mathrsfs}
\usetikzlibrary{arrows}



\begin{document}

    % 1. Foto de fondo
    % 2. Título
    % 3. Encabezado Izquierdo
    % 4. Color de fondo
    % 5. Coord x del titulo
    % 6. Coord y del titulo
    % 7. Fecha

    
    % 1. Foto de fondo
% 2. Título
% 3. Encabezado Izquierdo
% 4. Color de fondo
% 5. Coord x del titulo
% 6. Coord y del titulo
% 7. Fecha

\newcommand{\portada}[7]{

    \portadaBase{#1}{#2}{#3}{#4}{#5}{#6}{#7}
    \portadaBook{#1}{#2}{#3}{#4}{#5}{#6}{#7}
}

\newcommand{\portadaExamen}[7]{

    \portadaBase{#1}{#2}{#3}{#4}{#5}{#6}{#7}
    \portadaArticle{#1}{#2}{#3}{#4}{#5}{#6}{#7}
}




\newcommand{\portadaBase}[7]{

    % Tiene la portada principal y la licencia Creative Commons
    
    % 1. Foto de fondo
    % 2. Título
    % 3. Encabezado Izquierdo
    % 4. Color de fondo
    % 5. Coord x del titulo
    % 6. Coord y del titulo
    % 7. Fecha
    
    
    \thispagestyle{empty}               % Sin encabezado ni pie de página
    \newgeometry{margin=0cm}        % Márgenes nulos para la primera página
    
    
    % Encabezado
    \fancyhead[L]{\helv #3}
    \fancyhead[R]{\helv \nouppercase{\leftmark}}
    
    
    \pagecolor{#4}        % Color de fondo para la portada
    
    \begin{figure}[p]
        \centering
        \transparent{0.3}           % Opacidad del 30% para la imagen
        
        \includegraphics[width=\paperwidth, keepaspectratio]{assets/#1}
    
        \begin{tikzpicture}[remember picture, overlay]
            \node[anchor=north west, text=white, opacity=1, font=\fontsize{60}{90}\selectfont\bfseries\sffamily, align=left] at (#5, #6) {#2};
            
            \node[anchor=south east, text=white, opacity=1, font=\fontsize{12}{18}\selectfont\sffamily, align=right] at (9.7, 3) {\textbf{\href{https://losdeldgiim.github.io/}{Los Del DGIIM}}};
            
            \node[anchor=south east, text=white, opacity=1, font=\fontsize{12}{15}\selectfont\sffamily, align=right] at (9.7, 1.8) {Doble Grado en Ingeniería Informática y Matemáticas\\Universidad de Granada};
        \end{tikzpicture}
    \end{figure}
    
    
    \restoregeometry        % Restaurar márgenes normales para las páginas subsiguientes
    \pagecolor{white}       % Restaurar el color de página
    
    
    \newpage
    \thispagestyle{empty}               % Sin encabezado ni pie de página
    \begin{tikzpicture}[remember picture, overlay]
        \node[anchor=south west, inner sep=3cm] at (current page.south west) {
            \begin{minipage}{0.5\paperwidth}
                \href{https://creativecommons.org/licenses/by-nc-nd/4.0/}{
                    \includegraphics[height=2cm]{assets/Licencia.png}
                }\vspace{1cm}\\
                Esta obra está bajo una
                \href{https://creativecommons.org/licenses/by-nc-nd/4.0/}{
                    Licencia Creative Commons Atribución-NoComercial-SinDerivadas 4.0 Internacional (CC BY-NC-ND 4.0).
                }\\
    
                Eres libre de compartir y redistribuir el contenido de esta obra en cualquier medio o formato, siempre y cuando des el crédito adecuado a los autores originales y no persigas fines comerciales. 
            \end{minipage}
        };
    \end{tikzpicture}
    
    
    
    % 1. Foto de fondo
    % 2. Título
    % 3. Encabezado Izquierdo
    % 4. Color de fondo
    % 5. Coord x del titulo
    % 6. Coord y del titulo
    % 7. Fecha


}


\newcommand{\portadaBook}[7]{

    % 1. Foto de fondo
    % 2. Título
    % 3. Encabezado Izquierdo
    % 4. Color de fondo
    % 5. Coord x del titulo
    % 6. Coord y del titulo
    % 7. Fecha

    % Personaliza el formato del título
    \pretitle{\begin{center}\bfseries\fontsize{42}{56}\selectfont}
    \posttitle{\par\end{center}\vspace{2em}}
    
    % Personaliza el formato del autor
    \preauthor{\begin{center}\Large}
    \postauthor{\par\end{center}\vfill}
    
    % Personaliza el formato de la fecha
    \predate{\begin{center}\huge}
    \postdate{\par\end{center}\vspace{2em}}
    
    \title{#2}
    \author{\href{https://losdeldgiim.github.io/}{Los Del DGIIM}}
    \date{Granada, #7}
    \maketitle
    
    \tableofcontents
}




\newcommand{\portadaArticle}[7]{

    % 1. Foto de fondo
    % 2. Título
    % 3. Encabezado Izquierdo
    % 4. Color de fondo
    % 5. Coord x del titulo
    % 6. Coord y del titulo
    % 7. Fecha

    % Personaliza el formato del título
    \pretitle{\begin{center}\bfseries\fontsize{42}{56}\selectfont}
    \posttitle{\par\end{center}\vspace{2em}}
    
    % Personaliza el formato del autor
    \preauthor{\begin{center}\Large}
    \postauthor{\par\end{center}\vspace{3em}}
    
    % Personaliza el formato de la fecha
    \predate{\begin{center}\huge}
    \postdate{\par\end{center}\vspace{5em}}
    
    \title{#2}
    \author{\href{https://losdeldgiim.github.io/}{Los Del DGIIM}}
    \date{Granada, #7}
    \thispagestyle{empty}               % Sin encabezado ni pie de página
    \maketitle
    \vfill
}
    \portadaExamen{ffccA4.jpg}{Geometría III\\Examen XIII}{Geometría III. Examen XIII}{MidnightBlue}{-8}{28}{2023-2024}{Jesús Muñoz Velasco}

    
    \begin{description}
        \item[Asignatura] Geometría III.
        \item[Curso Académico] 2023-24.
        \item[Grado] Doble Grado en Ingeniería Informática y Matemáticas.
        \item[Grupo] Único.
        \item[Profesor] Antonio Martínez López.
        \item[Descripción] Convocatoria Extraordinaria.
        \item[Fecha] 17 de febrero de 2024.
        \item[Duración] 3 horas.
    
    \end{description}
    \newpage

    \begin{ejercicio}[2,5 puntos]
        Consideramos los subespacios afines de $\bb{R}^3$ dados por:
        \begin{gather*}
            S \equiv \{(x,y,z) \in \bb{R}^3 : x-y+z=2\}, \hspace*{1cm} T=(0,-1,0)+\cc{L}\{(1,1,1)\}.
        \end{gather*}

        Comprobar que $S$ y $T$ son suplementarios (o complementarios) afines. Calcula $\pi_{S,T}$ y $\sigma_{S,T}$ sobre $S$ en la dirección de $T$, dando sus ecuaciones matriciales respecto de $\cc{R}_0\equiv\{(0,0,0), \cc{B}_0\}$
        \\ \\
        Para comprobar que son suplementarios, deberemos comprobar que tienen intersección no vacia y que $\bb{R}^3 = \vec{S} \oplus \vec{T}$. Para ello calculamos en primer lugar unas ecuaciones cartesianas de $T$ en $\cc{R}_0$:
        \begin{equation*}
            \begin{vmatrix}
                1 & x \\ 
                1 & y + 1 \\
                1 & z
            \end{vmatrix}
        \end{equation*}
        Cuyo rango tiene que ser uno, es decir:
        \begin{align*}
            \begin{vmatrix}
                1 & x \\ 
                1 & y + 1 \\
            \end{vmatrix} = 0 \Rightarrow y + 1 = x
            \quad &
            \begin{vmatrix}
                1 & x \\ 
                1 & z \\
            \end{vmatrix} = 0 \Rightarrow z = x
        \end{align*}
        Como estamos trabajando en el mismo sistema de referencia $\cc{R}_0$, tenemos que:
        \begin{equation*}
            R \cap S = \begin{cases}
                z = x \\ 
                y + 1 = x \\
                x - y + z = 2
            \end{cases}
        \end{equation*}
        Que es un SCD con solución $S \cap T =(1,0,1)$ . Usando la fórmula de las dimensiones se obtiene además que 
        su suma es $\bb{R}^3$. Vamos a calcular ahora la proyección afin de $S$ paralela a $T$. Como se ha visto en teoria, dado $\{p_0\} = S \cap T$:
        \begin{equation*}
            \pi_{S,T} (p) = p_0 + \pi_{\vec{S},\vec{T}}(\vec{p_0 p})
        \end{equation*}
        Como $\vec{p_0p} = (x -1, y, z - 1)$, % Completar
    \end{ejercicio}

    \begin{ejercicio}[2,5 puntos]
        Razona en cada caso qué movimiento rígido resulta de las siguientes composiciones:

        \begin{enumerate}
            \item (1,25 Puntos) La composición en el espacio de un giro con una traslación.
                \\ \\
                Sea $f = G_{\theta, L} \circ t_v$, donde $L$ una recta, $\theta \in ]0, \pi[$ y $v \in \bb{R}^3 \backslash \{0\}$. Entonces $|\vec{f}| = 1$ por lo que se trata de una isometría directa. Distinguimos en función del valor de $v$:
                \begin{itemize}
                    \item Si $v \in \vec{L}$, se trata de un giro con deslizamiento, es decir, un movimiento helicoidal
                    \item Si $v \notin \vec{L}$, veamos que se trata de un giro:
                        \begin{equation*}
                            \vec{f} = \vec{G_{\theta, L}} \circ \vec{t_v} = G_{\theta, \vec{L}} \circ Id_{\bb{R}^3} = G_{\theta, \vec{L}}
                        \end{equation*}
                \end{itemize}
            \item (1,25 Puntos) La composición en el plano de un giro y una simetría con desplazamiento.
                \\ \\
                Sea $f = G_{\theta, p} \circ \sigma_R \circ t_v$ con $p \in \bb{R}^2, \theta \in ]0, \pi[$, $R$ una recta y $v \in \bb{R}^2 \backslash \{0\}$. Entonces $|\vec{f}| = -1$ por lo que se trata de una isometría inversa. Como la composición de un giro con una simetría es una simetría axial, llamando a esta $\sigma_S$, tenemos que:
                f = $\sigma_S \circ t_v$. Por lo que se trata de una simetría con o sin desplazamiento.
        \end{enumerate}

    \end{ejercicio}

    \begin{ejercicio}[2,5 puntos]
        Enuncia y demuestra el teorema de Desargues en el espacio afín y en el proyectivo de dimensión 2.
    \end{ejercicio}

    \begin{ejercicio}[2,5 puntos]
        Clasifica euclídeamente la siguiente cónica buscando el sistema de referencia en el que adopta su ecuación reducida. Calcula además sus elementos euclídeos:
        \begin{gather*}
            -16-8x+x^2+8y+2xy+y^2=0
        \end{gather*}
        Tenemos que la matriz asociada a dicha conica es la siguiente:
        \begin{equation*}
            \hat{A} = \left(
                \begin{array}{c|cc}
                    -16 & -4 & 4  \\
                    \hline
                    -4 & 1 & 1  \\
                    4 & 1 & 1  \\
                \end{array}
            \right) = \left(\begin{array}{c|c}
                a & z  \\
                \hline 
                z^t & A 
            \end{array} \right)
        \end{equation*}
        Diagonalizamos $A$, su polinomio característico es:
        \begin{equation*}
            \lambda^2 -2\lambda = 0 \Leftrightarrow \begin{cases}
                \lambda_1 = 0 \\
                \lambda_2 = 2
            \end{cases}
        \end{equation*}
        Sus vectores propios asociados son:
        \begin{align*}
            V_2 = \left\{ (x,y) \in \bb{R}^2 \, | \,  \begin{pmatrix} 
                    -1 & 1 \\ 
                    1 & -1
                    \end{pmatrix} \begin{pmatrix} x \\ y \end{pmatrix} = 0 \right\} = \cc{L} \left\{ \frac{1}{\sqrt{2}} (1, 1) \right\} \\
                    \\
            V_0 = \left\{ (x,y) \in \bb{R}^2 \, | \,  \begin{pmatrix} 
                    1 & 1 \\ 
                    1 & 1
                    \end{pmatrix} \begin{pmatrix} x \\ y \end{pmatrix} = 0 \right\} = \cc{L} \left\{ \frac{1}{\sqrt{2}} (1, -1) \right\} \\
        \end{align*}
        Por lo que se trata de una parábola. El eje entonces de la parábola es:
        \begin{align*}
            e = \left\{ (x,y) \in \bb{R}^2 \, | (0,1,1) \hat{A} (1, x, y) = 0 \right\} = \cc{L} \left\{ (x,y)^t \in \bb{R}^2  \, | \, x = -y\right\} = (0,0) + \cc{L} \left\{ (1, -1) \right\}
        \end{align*}
        Para calcular el vértice de la parábola, tenemos que resolver el siguiente sistema de ecuaciones:
        \[
        \left\{
        \begin{array}{l}
            x + y = 0 \\
            -16 - 8x + x^2 + 8y + 2xy + y^2 = 0
        \end{array}
    \right\}
        \Rightarrow v = (-1, 1)
        \]
        Sea pues el sistema de referencia $\cc{R}^{'} = \{ v, B = \{ \frac{1}{\sqrt{2}}(1,1), \frac{1}{\sqrt{2}}(1,-1) \} \}$. En $\cc{R}^{'}$ su matriz asociada queda:
        \begin{equation*}
            \left(
                \begin{array}{c|cc}
                    0 & 0 & \lambda  \\
                    \hline
                    0 & 2 & 0  \\
                    \lambda & 0 & 0  \\
                \end{array}
            \right) 
        \end{equation*}
        Para calcular $\lambda$ sabemos que el determinante es invariante por semejanzas por tanto:
        \begin{equation*}
            |\hat{A}| = -64 = -2\lambda^2 \Rightarrow \lambda = \pm\sqrt{32}
        \end{equation*}
    \end{ejercicio}
    Por tanto su ecuación en dicho sistema es $2x^2 = -2\lambda y $. Como el coeficiente de x es positivo, se tiene la solución negativa de $\lambda$. Si la distancia del vertice al foco es $c$, entonces:
    % Completar 
\end{document}
